\section{Conclusion}
\label{sec:conclusion}

We have introduced \ecagbench{}, a benchmark of 332 constructive problems in algebraic geometry designed to evaluate a mode of mathematical reasoning---object construction---that is central to mathematical practice but underrepresented in existing LLM benchmarks.

\paragraph{Key findings.}
Our pilot evaluation of Claude Sonnet~4.6 on a stratified 100-problem subset reveals:
\begin{enumerate}[nosep]
  \item \textbf{A steep difficulty gradient:} The model achieves 78\% accuracy overall, but this masks a 68-percentage-point drop from Level~1 (95\%) to Level~3 (27\%), confirming that \ecagbench{} effectively discriminates between routine constructions and deep problems requiring synthesis of advanced techniques.
  \item \textbf{Counterexamples are harder than examples:} The 15-percentage-point gap (80\% vs.\ 65\%) confirms that understanding theorem boundaries requires deeper reasoning than applying theorems directly.
  \item \textbf{Timeouts as a failure mode:} 18\% of problems caused the model to time out (>600s), with the timeout rate reaching 45.5\% at Level~3, suggesting that the hardest problems cause the model to fail silently before producing any response.
  \item \textbf{Chapter-dependent performance:} Accuracy ranges from 93\% (Cohomology) to 65\% (Moduli), reflecting the varying difficulty of constructive reasoning across algebraic geometry subfields.
  \item \textbf{Failure modes are structured:} The six failure modes (F1--F6) provide a systematic taxonomy for characterizing model weaknesses, with property confusion (F2) and incomplete verification (F6) observed in the pilot.
\end{enumerate}

\paragraph{Limitations.}
Our evaluation has several limitations that should be addressed in future work:
\begin{itemize}[nosep]
  \item \textbf{Single model:} We evaluate only Claude Sonnet~4.6. Multi-model comparison (GPT-4o, Gemini 2.5 Pro, DeepSeek-R1~\citep{deepseek2025r1}, Llemma~\citep{azerbayev2024llemma}) is needed to determine whether the observed patterns are model-specific.
  \item \textbf{Self-evaluation bias:} Both the problem solver and the automated scorer are LLMs from the same model family, introducing potential concordance bias. Among the 82 problems that received responses, 95.1\% were scored as correct, which may reflect systematic overrating. Human expert scoring is needed to validate these results (see Section~\ref{sec:self-eval-limitations}).
  \item \textbf{Single-turn only:} Multi-turn interaction, where the model can refine constructions based on CAS feedback, may significantly improve performance and is an important evaluation modality to explore.
  \item \textbf{Domain specificity:} Results in algebraic geometry may not generalize to other mathematical domains, though we expect similar example/counterexample patterns.
  \item \textbf{CAS coverage:} While 93.1\% of problems use CAS verification, the CAS cannot verify all properties, and some verifications require specialized scripts that have not yet been fully implemented.
\end{itemize}

\paragraph{Future work.}
We plan several extensions:
\begin{enumerate}[nosep]
  \item \textbf{Full-scale evaluation with human scoring:} Evaluate multiple frontier models on the complete 332-problem benchmark with human expert scoring to establish unbiased baselines.
  \item \textbf{Multi-turn interaction:} Evaluate models in an interactive setting where they receive CAS feedback and can revise their constructions.
  \item \textbf{Lean~4 integration:} Formalize a subset of problems in Lean~4, enabling formal verification as an additional evaluation layer.
  \item \textbf{Fine-tuning studies:} Investigate whether fine-tuning on algebraic geometry textbooks~\citep{hartshorne1977algebraic, vakil2017rising, eisenbud1995commutative} or the accompanying \ecagbench{} textbook improves constructive reasoning.
  \item \textbf{Community contributions:} The open-source nature of the project (CC-BY-4.0 license) invites community contributions of new problems, verification scripts, and evaluation results.
  \item \textbf{Cross-domain extension:} Develop analogous benchmarks for arithmetic geometry, algebraic topology, and commutative algebra to investigate whether the example/counterexample performance gap generalizes.
\end{enumerate}

\paragraph{Broader impact.}
Constructive mathematical reasoning is a prerequisite for using LLMs as research tools in mathematics.
If models cannot construct examples and counterexamples, they cannot assist mathematicians in the exploratory phase of research---the phase where intuition is built and conjectures are tested.
The pronounced difficulty of counterexample construction is particularly significant: as Lakatos~\citep{lakatos1976proofs} argued, mathematical progress depends on the ability to find refutations, and this ability appears to be a current blind spot for language models.
\ecagbench{} provides a concrete measure of progress toward closing this gap.
